\documentclass[12pt]{amsart} 

\input{./latex-preamble/cm-preamble.tex}

%\graphicspath{ {./diagrams/} }

\newtheorem{lemma}{Lemma}

\newcommand{\D}{{\mathcal D}}

\begin{document}

{Hartshorne} \hfill {4-23-2025}

\bigskip\bigskip

\begin{center}

{\sc Chapter III.5, Week 12}

\end{center}

\begin{center}

  {Noah Parker} 
    
\end{center}

\bigskip\bigskip

\begin{enumerate}
  \item[5.1.] {
      Let $X$ be a projective scheme over a field $k$, and let $\F$ be a coherent sheaf on $X$.
      We define the \emph{Euler characteristic} of $\F$ by
      \[
        \chi(\F)=\sum_i(-1)^i\dim_kH^i(X,\F).
      \] 
      If
      \[
        0\to \F'\to \F\to \F''\to 0
      \] 
      is a short exact sequence of coherent sheaves on $X$, show that $\chi(\F)=\chi(\F')+\chi(\F'')$.
    }
\end{enumerate}

\begin{proof}
  First, for an exact sequence of vector spaces $0\to V_0\xrightarrow{f_0}\cdots\xrightarrow{f_{n-1}} V_n\xrightarrow{f_n}0$, consider $\sum_i(-1)^i\dim V_i$.
  We have 
  \begin{align*}
    \sum_{i=0}^n(-1)^i\dim V_i 
    &= \sum_{i=1}^n(-1)^i\dim \ker f_i +\sum_{i=0}^{n-1}(-1)^i\dim \im f_i \\
    &= \sum_{i=1}^n(-1)^i \dim \ker f_i +\sum_{i=0}^{n-1} (-1)^i\dim\ker f_{i+1} \\
    &=0.
  \end{align*}
  Thus, the long exact sequence associated to $0\to \F'\to \F\to \F''\to 0$ gives us
  \begin{align*}
    \chi(\F')-\chi(\F)+\chi(\F'')= 0,
  \end{align*}
  as desired.
\end{proof}

\begin{enumerate}
  \item[5.2.] {
      \begin{enumerate}
        \item Let $X$ be a projective scheme over a field $k$, let $\O_X(1)$ be a very ample invertible sheaf on $X$ over $k$, and let $\F$ be a coherent sheaf on $X$.
          Show that there is a polynomial $P(z)\in \QQ[z]$, such that $\chi(\F(n))=P(n)$ for all $n\in \ZZ$.
          We call $P$ the \emph{Hilbert polynomial} of $\F$ with respect to the sheaf $\O_X(1)$.

          [\emph{Hint}: Use induction on $\dim \Supp \F$, general properties of numerical polynomials (I, 7.3), and suitable exact sequences
          \[
            0\to\R\to \F(-1)\to \F\to \mathscr{D}\to 0.]
          \] 
        \item Now let $X=\PP^r_k$, and let $M=\Gamma_*(\F)$, considered as a graded $S=k[x_0,\ldots,x_r]$-module.
          Use (5.2) to show that the Hilbert polynomial of $\F$ just defined is the same as the Hilbert polynomial of $M$ defined in (I, 7).
      \end{enumerate}
    }
\end{enumerate}

\begin{proof}
  (a)

  (b) Let $P(z)$ be the Hilbert poylnomial of $\F$ defined in (a).
  It suffices to show that $P(n)=\dim M_n =\dim_k H^0(X,\F(n))$ for $n\gg 0$.
  However, (5.2) gives us
  \begin{align*}
    P(n)
    &=\chi(\F(n)) \\
    &=\sum_{i=0}^n (-1)^iH^i(X,\F(n)) \\
    &=H^0(X,\F(n))
  \end{align*}
  for $n\gg 0$, so we're done.
\end{proof}

\begin{enumerate}
  \item[5.3.] {
      \emph{Arithmetic Genus.}
      Let $X$ be a projective scheme of dimension $r$ over a field $k$.
      We define the \emph{arithmetic genus} $p_a$ of $X$ by
      \[
        p_a(X)= (-1)^r (\chi(\O_X)-1).
      \] 
      Note that it depends only on $X$, not on any projective embedding.
      \begin{enumerate}
        \item If $X$ is integral, and $k$ algebraically closed, show that $H^0(X,\O_X)\cong k$, so that
          \begin{align*}
            p_a(X)=\sum_{i=0}^{r-1}(-1)^i\dim_k H^{r-i}(X,\O_X).
          \end{align*}
          In particular, if $X$ is a curve, we have
          \[
            p_a(X)=\dim_k H^1(X,\O_X).
          \] 
          [\emph{Hint}: Use (I, 3.4).]
        \item If $X$ is a closed subvariety of $\PP^r_k$, show that this $p_a(X)$ coincidese with the one defined in (I, Ex. 7.2), which apparently depended on the projective embedding.
        \item If $X$ is a nonsingular projective curve over an algebraically closed field $k$, show that $p_a(X)$ is in fact a \emph{birational} invariant.
          Conclude that a nonsingular plane curve of degree $d\geqs 3$ is not rational.
          (This gives another proof of (II, 8.20.3) where we used the geometric genus.)
      \end{enumerate}
    }
\end{enumerate}

\begin{proof}
  (a) We have that $H^0(X,\O_X)\cong k$ by (I, 3.4).
  Then
  \begin{align*}
    p_a(X)
    &= (-1)^r (\chi(\O_X)-1) \\
    &= (-1)^r (\chi(\O_X)- \dim_k H^0(X,\O_X)) \\
    &= (-1)^r (\sum_{i=1}^r (-1)^i\dim_k H^i(X,\O_X)) \\
    &= \sum_{i=1}^r (-1)^{r-i}\dim_kH^i(X,\O_X) \\
    &=\sum_{i=0}^{r-1}(-1)^i \dim_k H^i(X,\O_X).
  \end{align*}
  If $X$ is a curve, then $r=1$, so the result holds.

  (b)

  (c)
\end{proof}

\begin{enumerate}
  \item[5.4.] {
      Recall from (II, Ex. 6.10) the definition of the Grothendieck group $K(X)$ of a noetherian scheme $X$.
      \begin{enumerate}
        \item Let $X$ be a projective scheme over a field $k$, and let $\O_X(a)$ be a very ample invertible sheaf on $X$.
          Show that there is a (unique) additive homomorphism
          \[
            P:K(X)\to \QQ[z]
          \] 
          such that for each coherent sheaf $\F$ on $X$, $P(\gamma(\F))$ is the Hilbert polynomial of $\F$ (Ex. 5.2).
        \item Now let $X=\PP^r_k$.
          For each $i=0,1,\ldots, r$, let $L_i$ be a linear space of dimension $i$ in $X$.
          Then show that
          \begin{enumerate}
            \item $K(X)$ is the free abelian group generated by $\{\gamma(\O_{L_i})\mid i=0,\ldots, r\}$, and
            \item the map $P:K(X)\to \QQ[z]$ is injective.
          \end{enumerate}
          [\emph{Hint}: Show that (i)$\implies$(ii).
          Then prove (i) and (ii) simultaneously, by induction on $r$, using (II, Ex. 6.10c).]
      \end{enumerate}
    }
\end{enumerate}

\begin{proof}
  (a) The obvious map on the free abelian group generated by all coherent sheaves of $X$ descends to $K(X)$ by (Ex. 5.1) above.

  (b) If (i) holds, then $P$ is injective because $P(\gamma(\O_{L_i}))$
\end{proof}

\begin{enumerate}
  \item[5.7.] {
      Let $X$, $Y$ be proper schemes over a noetherian ring $A$.
      We denote by $\L$ an invertible sheaf.
      \begin{enumerate}
        \item  If $\L$ is ample on $X$, and $Y$ is any closed subscheme of $X$, then $i^*\L$ is ample on $Y$, where $i:Y\to X$ is the inclusion.
        \item $\L$ is ample on $X$ if and only if $\L_\red =\L\otimes \O_{X_\red}$ is ample on $X_\red$.
        \item Suppose $X$ is reduced.
          Then $\L$ is ample on $X$ if and only if $\L\otimes \O_{X_i}$ is ample on $X_i$ for each irreducible component $X_i$ of $X$.
        \item Let $f:X\to Y$ be a finite surjective morphism, and let $\L$ be an invertible sheaf on $Y$.
          Then $\L$ is ample on $Y$ if and only if $f^*\L$ is ample on $X$.
      \end{enumerate}
      [\emph{Hints}: Use (5.3) and compare (Ex. 3.1, Ex. 3.2, Ex. 4.1, Ex. 4.2).]
    }
\end{enumerate}

\begin{proof}
  (a) Let $\F$ be a coherent sheaf on $Y$.
  $i^*$ is a left adjoint functor and tensor products are direct limits, so $(i^*\L)^n=i^*(\L^n)$.
  Thus, we may omit the parentheses without ambiguity.
  By (II, Ex. 5.1d), we have
  \[
    i_*(\F\otimes_{\O_Y} i^*\L^n)\cong i_*\F\otimes_{\O_X} \L^n.
  \] 
  $X$, $Y$ are noetherian and $i$ is a closed immersion, so $i_*\F$ is coherent by (II, Ex. 5.5).
  Then (5.3) gives us $n_0>0$ such that, for all $n>0$ and $p>0$,
  \begin{align*}
    H^p(Y,\F\otimes i^*\L^n)
    &\cong H^p(X,i_*\F\otimes \L^n)
    \cong 0
  \end{align*}
  by (2.10). 
  By (5.3) once more, we conclude that $i^*\L$ is ample.

  (b) We begin by noting that the immersion $i:X_\red\inj X$ is a bijection, so $i^*\L = \L\otimes \O_X=\L_\red$.
  If $\L$ is ample on $X$, then $\L_\red$ is ample on $X_\red$ by (a).
  Suppose then that $\L_\red$ is ample on $X_\red$ and fix a coherent sheaf $\F$ on $X$.
  As in (3.1), noetherianity of $X$ gives a finite filtration $\F\supset\N\F\supset \cdots \supset \N^n\F=0$, and, thus, the exact sequences
  \[
    0\to \N^{k+1}\F\to\N^k\F\to \frac{\N^k\F}{\N^{k+1}\F}\to 0
  \]
  for $0\leqs k<n$.

  Note that the quotient sheaves are coherent by coherence of $\N^k\F$ for all $k\geqs 0$.
  Then
  \[
    \N\subset \Ann \left(\frac{\N^k\F}{\N^{k+1}\F}\right),
  \] 
  supplies us with sheaves $\G_k$ on $X_\red$ such that $i_*\G_k$.
  In fact, the $\G_k$ are coherent.
  This is because $\im \ol f=\im (\ol f\circ \pi)=\im f$ whenever $\pi$ is surjective, so in particular when $\pi:A\to A/\a$ is the projection.
  Then (5.3) supplies $n_0>0$ such that, for all $n\geqs n_0$ and $p>0$, $H^p(X_\red, \G_k\otimes \L^n)\cong 0$.

  Applying the projection formula (II, 5.1) gives us
  \begin{align*}
    H^\bullet(X,\frac{\N^k\F}{\N^{k+1}\F}\otimes \L^n)
    &\cong H^\bullet(X_\red, \G_k\otimes \L^n).
  \end{align*}
  Tensoring our original exact sequence with $\L^n$ for $n\geqs n_0$ and taking the long exact sequence gives us exactness of
  \begin{align*}
    H^p(X,\N^{k+1}\F\otimes \L^n)\to H^p(X,\N^k\F\otimes \L^n)\to H^p(X_\red,\G_k\otimes \L^n) =0
  \end{align*}
  for $p>0$.
  $\N^n\F=0$, so it follows by induction on $n-k$ that
  \[
    H^p(X,\F\otimes \L^n)=0
  \] 
  for all $p>0$, whence $\L$ is ample by (5.3).

  (c) $\L_i:=\L\otimes X_i$ is slightly abusive notation for the inverse image of the inclusion $X_i\hookrightarrow X$, where $X_i$ is given the reduced induced subscheme structure.
  Thus, one direction is a consequence of (a).
  It suffices to show that $\L$ is ample if $\L_i$ is ample for each irreducible component $X_i$ of $X$.

  Let $X=X_1\cup\cdots\cup X_r$ be an irreducible decomposition of $X$, finite since $X$ is noetherian.
  Suppose that $\L_i$ is ample for each $i=1,\ldots, r$, and consider the finite filtration $\F\supset \I_{X_1}\F\supset\cdots \supset \I_{X_1\cup\cdots\cup X_r}\F=0$, where each $X_1\cup\cdots \cup X_k$ is given the reduced induced subscheme structure.
  We once again receive exact sequences
  \[
    0\to \I_{Y_1\cup\cdots\cup Y_{k+1}}\F\to \I_{Y_1\cup\cdots \cup Y_k}\F\to \frac{\I_{Y_1\cup\cdots \cup Y_k}\F} {\I_{Y_1\cup\cdots \cup Y_{k+1}} \F} \to 0
  \]
  for all $0\leqs k<r$, where
  \[
    \I_{Y_{k+1}}\subset \Ann \left(\frac{\I_{Y_1\cup\cdots \cup Y_k}\F} {\I_{Y_1\cup\cdots \cup Y_{k+1}} \F}\right).
  \]
  As above, we have coherent sheaves $\G_{k+1}$ on $Y_{k+1}$ such that
  \[
    H^\bullet(X,\frac{\I_{Y_1\cup\cdots \cup Y_k}\F} {\I_{Y_1\cup\cdots \cup Y_{k+1}} \F})\cong H^\bullet(Y_{k+1},\G_{k+1}).
  \]
  Then (5.3) supplies us with $n_0>0$ such that
  \[
    H^p(Y_{k+1},\G_{k+1}\otimes \L_{k+1}^n) = 0
  \] 
  for all $p>0$ and $n\geqs n_0$.
  Once again, this gives us surjectivity of the maps
  \[
    H^p(X,\I_{Y_1\cup\cdots \cup Y_{k+1}}\F)\to H^p(X,\I_{Y_1\cup\cdots\cup Y_k}\F)
  \] 
  whenever $p>0$, so the result follows by induction on $r-k$.

  (d) Suppose first that $\L$ is ample on $X$.
  Let $\F$ be a coherent sheaf on $Y$.
  By (II, Ex. 5.1d), we have
  \[
    f_*(\F\otimes_{\O_Y} f^*\L^n)\cong f_*\F\otimes_{\O_X} \L^n.
  \] 
  $X$, $Y$ are noetherian and $f$ is a finite map, so $f_*\F$ is coherent by (II, Ex. 5.5).
  Then (5.3) gives us $n_0>0$ such that, for all $n>0$ and $p>0$,
  \begin{align*}
    H^p(Y,\F\otimes f^*\L^n)
    &\cong H^p(X,f_*\F\otimes \L^n)
    \cong 0
  \end{align*}
  by (2.10). 
  By (5.3) once more, we conclude that $f^*\L$ is ample.

  Suppose now that $f^*\L$ is ample.
\end{proof}

\begin{enumerate}
  \item[5.8.] {
      Prove that every one-dimensional proper scheme $X$ over an algebraically closed field $k$ is projective.
      \begin{enumerate}
        \item If $X$ is irreducible and nonsingular, then $X$ is projective by (II, 6.7).
        \item If $X$ is integral, let $\wtilde X$ be its normalization (II, Ex. 3.8).
          Show that $\wtilde X$ is complete and nonsingular, hence projective by (a).

          Let $f:\wtilde X\to X$ be the projection.
          Let $\L$ be a very ample invertible sheaf on $\wtilde X$.
          Show there is an effective divisor $D=\sum P_i$ on $\wtilde X$ with $\L(D)\cong \L$, and such that $f(P_i)$ is a nonsingular point of $X$, for each $i$.

          Conclude that there is an invertible sheaf $\L_0$ on $X$ with $f^*\L_0\cong \L$.
          Then use (Ex. 5.7d), (II, 7.6), and (II, 5.16.1) to show that $X$ is projective.
        \item If $X$ is reduced, but not necessarily irreducible, let $X_1,\ldots, X_r$ be the irreduciblee components of $X$.
          Use (Ex. 4.5) to show that $\Pic X\to \bigoplus \Pic X_i$ is surjective.
          Then use (Ex. 5.7c) to show $X$ is projective.
        \item Finally, if $X$ is any one dimensional proper scheme over $k$, use (2.7) and (Ex. 4.6) to show that $\Pic X\to \Pic X_\red$ is surjective.
          Then use (Ex. 5.7b) to show $X$ is projective.
      \end{enumerate}
    }
\end{enumerate}

\begin{proof}
  (a) This is exact (ii)$\implies$(i) of (II, 6.7).

  (b) 
\end{proof}

\end{document}
